\section{Survey Instrument}

\subsection{Vignette Experiment and Profile Construction}

\noindent\textbf{Instructions to respondents.} In the following part of the survey, you will see descriptions of several people who
have come to Germany. The people described differ in several characteristics. Please
read each description carefully and answer the questions that follow based on what you
think is most likely for the person described. There are no right or wrong answers; we
are interested in your beliefs about each person.\medskip

\noindent\textbf{Illustrative profiles.} The following profiles illustrate how the randomized attribute combinations would
appear to respondents.

\medskip

\noindent\textit{Example 1.}

\noindent
A male from Syria came to Germany. He came to Germany as a refugee after fleeing
the war in his country. He has no formal vocational qualification, speaks German only
to a limited extent, and is currently looking for stable employment.

\medskip

\noindent\textit{Example 2.}

\noindent
A female from Ukraine came to Germany. She came to Germany as an economic migrant
in search of better employment and economic opportunities in Germany. She has
completed a university degree, speaks German well, and is currently looking for stable
employment.

\medskip

\noindent\textit{Example 3.}

\noindent
A female from Syria came to Germany. She came to Germany as an economic migrant
in search of better employment and economic opportunities in Germany. She has
completed a university degree, speaks German only to a limited extent, and is currently
looking for stable employment.

\medskip

\noindent\textit{Example 4.}

\noindent
A male from Ukraine came to Germany. He came to Germany as a refugee after fleeing
the war in his country. He has no formal vocational qualification, speaks German well,
and is currently looking for stable employment.


\subsection{Outcome Questions and Response Scales}
\label{app:outcomes}

After each vignette profile, respondents would answer the following five questions.
The same questions and response scales would be presented after each profile.

\medskip

\noindent\textbf{Q1.}
What do you think is the probability that this person will have a regular paid job
subject to social security contributions within the next two years?

\smallskip

\noindent
\textbf{Response:} 0--100\%, where 0\% means ``no chance'' and 100\% means
``certain.''

\medskip

\noindent\textbf{Q2.}
What do you think this person's monthly net income will be five years from now?

\smallskip

\noindent
\textbf{Response options:}
\begin{enumerate}
    \item Less than €1,500
    \item €1,500--€1,999
    \item €2,000--€2,499
    \item €2,500--€2,999
    \item €3,000--€3,499
    \item €3,500--€3,999
    \item €4,000 or more
\end{enumerate}

\medskip

\noindent\textbf{Q3.}
How likely do you think it is that this person would be treated less favorably
than an otherwise comparable German applicant when applying for a job?

\smallskip

\noindent
\textbf{Response:} 1--7, where 1 means ``very unlikely'' and 7 means
``very likely.''

\medskip

\noindent\textbf{Q4.}
How likely do you think it is that this person would face unequal treatment in the workplace, 
for example in promotion opportunities, wages, or interactions with colleagues? 

\smallskip

\noindent
\textbf{Response:} 1--7, where 1 means ``very unlikely'' and 7 means
``very likely.''

\medskip

\noindent\textbf{Q5.}
How likely do you think it is that this person will have regular social contacts
with Germans within the next five years? 

\smallskip

\noindent
\textbf{Response:} 1--7, where 1 means ``very unlikely'' and 7 means
``very likely.''

\subsection{Personal Contact and Experience Questions}

The following questions would be asked separately for Ukrainian and Syrian
migrants. The order of the two migrant groups would be randomized where feasible.

\medskip

\noindent\textbf{Q6.}
During the past 12 months, how often have you personally interacted with
Ukrainian/Syrian colleagues or other Ukrainian/Syrian people in connection
with your work?

\smallskip

\noindent\textbf{Response options:}

\begin{enumerate}
    \item Never
    \item Less than once a month
    \item About 1--3 times per month
    \item About once per week
    \item Several times per week
    \item About once per day
    \item Several times per day
    \item Not applicable / I have not worked during the past 12 months
\end{enumerate}

\noindent
\textit{Display/coding note:} ``Not applicable / I have not worked during the
past 12 months'' is treated as not applicable and coded as missing.

\medskip

\noindent\textbf{Q7.}
During the past 12 months, how often have you personally interacted with
Ukrainians/Syrians outside of work, for example as friends, neighbours,
acquaintances, through clubs, or in other social settings?

\smallskip

\noindent\textbf{Response options:}

\begin{enumerate}
    \item Never
    \item Less than once a month
    \item About 1--3 times per month
    \item About once per week
    \item Several times per week
    \item About once per day
    \item Several times per day
\end{enumerate}

\medskip

\noindent\textbf{Q8.}
Thinking about your personal interactions with Ukrainians/Syrians living in
Germany overall, how would you describe these experiences?

\smallskip

\noindent\textbf{Response:} 1--7, where 1 means ``Very negative,''
4 means ``Neither positive nor negative,'' and 7 means ``Very positive.''

\smallskip

\noindent
\textit{Display logic:} Q8 is displayed if the respondent does not answer
``Never'' to both Q6 and Q7.

\medskip

\noindent\textbf{Q9.}
Thinking about your personal interactions with Ukrainians/Syrians living in
Germany, does a particular experience immediately come to mind?

\smallskip

\noindent\textbf{Response options:}

\begin{itemize}
    \item Yes
    \item No
\end{itemize}

\smallskip

\noindent
\textit{Display logic:} Q9 is displayed if the respondent reports personal
contact with the respective migrant group in Q6 or Q7.

\medskip

\noindent\textbf{Q10.}
Thinking about this particular experience, how would you describe it?

\smallskip

\noindent\textbf{Response:} 1--7, where 1 means ``Very negative,''
4 means ``Neither positive nor negative,'' and 7 means ``Very positive.''

\smallskip

\noindent
\textit{Display logic:} Q10 is displayed only if the respondent answers
``Yes'' to Q9.

\subsection{Media Exposure and Information Questions}

The following questions would be asked separately for Ukrainian and Syrian
migrants. The order of the two groups would be randomized where feasible.

\medskip

\noindent\textbf{Q11.}
How often do you come across media information about Ukrainian/Syrian migrants
living in Germany through television, newspapers, online news, social media,
radio, podcasts, or other media?

\smallskip

\noindent\textbf{Response options:}

\begin{enumerate}
    \item Never
    \item Less than once a month
    \item About 1--3 times per month
    \item About once per week
    \item Several times per week
    \item About once per day
    \item Several times per day
\end{enumerate}

\medskip

\noindent\textbf{Q12.}
Overall, how would you describe the information you have encountered about
Ukrainian/Syrian migrants living in Germany?

\smallskip

\noindent
\textbf{Response:} 1--7, where 1 means ``Very negative,''
4 means ``Neither positive nor negative,'' and 7 means ``Very positive.''

\smallskip

\noindent
\textit{Display logic:} Q12 is displayed only if the respondent does not
answer ``Never'' to Q11.

\medskip

\noindent\textbf{Q13.}
When thinking about Ukrainian/Syrian migrants living in Germany, does a
particular piece of media information immediately come to mind?

\smallskip

\noindent\textbf{Response options:}

\begin{itemize}
    \item Yes
    \item No
\end{itemize}

\smallskip

\noindent
\textit{Display logic:} Q13 is displayed only if the respondent does not
answer ``Never'' to Q11.

\medskip

\noindent\textbf{Q14.}
Overall, how would you describe that piece of information?

\smallskip

\noindent
\textbf{Response:} 1--7, where 1 means ``Very negative,''
4 means ``Neither positive nor negative,'' and 7 means ``Very positive.''

\smallskip

\noindent
\textit{Display logic:} Q14 is displayed only if the respondent answers
``Yes'' to Q13.

\subsection{General-Belief Questions}

The following questions measure respondents' broader beliefs about migrant
integration and migrant communities in Germany.

\medskip

\noindent\textbf{Q15.}
Migrants tend to become better integrated into the German labor market the
longer they live in Germany.

\smallskip

\noindent
\textbf{Response:} 1--7, where 1 means ``Strongly disagree'' and
7 means ``Strongly agree.''

\medskip

\noindent\textbf{Q16.}
Migrants tend to become more socially integrated in Germany the longer they
live in the country.

\smallskip

\noindent
\textbf{Response:} 1--7, where 1 means ``Strongly disagree'' and
7 means ``Strongly agree.''

\medskip

\noindent
The following questions would be asked separately for Ukrainian and Syrian
migrants. The order of the two groups would be randomized where feasible.

\medskip

\noindent\textbf{Q17.}
The Ukrainian/Syrian community in Germany is well established.

\smallskip

\noindent
\textbf{Response:} 1--7, where 1 means ``Strongly disagree'' and
7 means ``Strongly agree.''

\medskip

\noindent\textbf{Q18.}
Ukrainian/Syrian communities in Germany help newly arrived
Ukrainians/Syrians find employment through personal networks.

\smallskip

\noindent
\textbf{Response:} 1--7, where 1 means ``Strongly disagree'' and
7 means ``Strongly agree.''

\medskip

\noindent\textbf{Q19.}
Being able to rely on an established Ukrainian/Syrian community reduces the
need for newly arrived Ukrainians/Syrians to interact regularly with Germans.

\smallskip

\noindent
\textbf{Response:} 1--7, where 1 means ``Strongly disagree'' and
7 means ``Strongly agree.''

\subsection{Respondent Characteristics}

The following questions collect respondent-level characteristics used in the
adjusted specifications. All displayed questions require a response before the respondent can proceed.
For sensitive respondent characteristics, ``Prefer not to say'' is provided
as an explicit response option. These responses are retained as explicit
categories in the raw survey data and coded as missing for the corresponding
variable in the empirical analysis.

\medskip

\noindent\textbf{Q20.}
What is your age?

\smallskip

\noindent
\textbf{Response:} Age in completed years.

\medskip

\noindent\textbf{Q21.}
What is your gender?

\smallskip

\noindent\textbf{Response options:}
\begin{itemize}
    \item Female
    \item Male
    \item Diverse / another gender
    \item Prefer not to say
\end{itemize}

\medskip

\noindent\textbf{Q22.}
What is the highest educational qualification you have completed?

\smallskip

\noindent\textbf{Response options:}
\begin{itemize}
    \item No formal school qualification
    \item Lower secondary school certificate (e.g., Hauptschule or equivalent)
    \item Intermediate secondary school certificate (e.g., Realschule or equivalent)
    \item Upper secondary school qualification (e.g., Abitur, Fachabitur, or equivalent)
    \item Completed vocational education or apprenticeship (Ausbildung)
    \item Bachelor's degree or equivalent
    \item Master's degree or equivalent
    \item Doctorate (PhD) or equivalent
    \item Other (please specify): \rule{3cm}{0.4pt}
    \item Prefer not to say
\end{itemize}

\medskip

\noindent\textbf{Q23.}
What is your current employment status?

\smallskip

\noindent\textbf{Response options:}
\begin{itemize}
    \item Employed full-time
    \item Employed part-time
    \item Self-employed
    \item Student / apprentice
    \item Unemployed and actively looking for work
    \item Not employed and not currently looking for work
    \item Retired
    \item Other
    \item Prefer not to say
\end{itemize}

\medskip

\noindent\textbf{Q24.}
What is your approximate monthly net household income (after taxes and social
security contributions)?

\smallskip

\noindent\textbf{Response options:}
\begin{itemize}
    \item Less than €1,500
    \item €1,500--€2,499
    \item €2,500--€3,499
    \item €3,500--€4,499
    \item €4,500--€5,499
    \item €5,500--€6,499
    \item €6,500--€7,499
    \item €7,500--€7,999
    \item €8,000 or more
    \item Prefer not to say
\end{itemize}

\medskip

\noindent\textbf{Q25.}
In which country were you born?

\smallskip

\noindent\textbf{Response options:}
\begin{itemize}
    \item Germany
    \item Other country: \rule{3cm}{0.4pt}
    \item Prefer not to say
\end{itemize}

\medskip

\noindent\textbf{Q26.}
In which country was your mother born?

\smallskip

\noindent\textbf{Response options:}
\begin{itemize}
    \item Germany
    \item Other country: \rule{3cm}{0.4pt}
    \item Don't know
    \item Prefer not to say
\end{itemize}

\medskip

\noindent\textbf{Q27.}
In which country was your father born?

\smallskip

\noindent\textbf{Response options:}
\begin{itemize}
    \item Germany
    \item Other country: \rule{3cm}{0.4pt}
    \item Don't know
    \item Prefer not to say
\end{itemize}

\medskip

\noindent\textbf{Q28.}
In which federal state (Bundesland) do you currently live?

\smallskip

\noindent
\textbf{Response:} Dropdown list containing all 16 German federal states.

\medskip

\noindent\textbf{Q29.}
What is your religious affiliation?

\smallskip

\noindent\textbf{Response options:}
\begin{itemize}
    \item Christian
    \item Muslim
    \item Jewish
    \item Other religious affiliation
    \item No religious affiliation
    \item Prefer not to say
\end{itemize}

\medskip

\noindent\textbf{Q30.}
In politics, people sometimes talk about ``left'' and ``right.'' Where would
you place yourself on a scale from 0 to 10?

\smallskip

\noindent
\textbf{Response:} 0--10, where 0 means ``Left,'' 5 means ``Centre,''
and 10 means ``Right.''

\medskip

\noindent\textbf{Coding and variable construction.} Age is recorded in completed years and can be included as a continuous
respondent-level variable. Gender, education, employment status, household
income, religion, and region are retained in their detailed categories in the
raw data and can be represented by categorical indicators in the empirical
analysis. Categories may be combined where necessary if individual categories
contain too few observations.

Migration background is constructed from Q25--Q27. The indicator
\(Mig_i\) equals 1 if the respondent or at least one parent was born outside
Germany, and 0 if the respondent and both parents were born in Germany.
``Don't know'' and ``Prefer not to say'' responses that prevent this
classification are treated as missing for this variable.

The respondent's federal state is retained in the raw data. Where relevant,
an indicator \(East_i\) can additionally be constructed to distinguish
respondents residing in eastern and western Germany. Political orientation,
\(Right_i\), is recorded on the original 0--10 scale, with higher values
indicating a more right-wing self-placement.
\medskip

\noindent\textbf{Attention Check.} The proposed survey would include one instructed-response attention check,
placed immediately after the vignette experiment and before the respondent
questionnaire. A possible item would be:

\medskip
\noindent
To show that you are paying attention, please select ``Somewhat agree'' below.

\smallskip
\noindent
\textit{Response options:} Strongly disagree; Disagree; Somewhat disagree;
Neither agree nor disagree; Somewhat agree; Agree; Strongly agree.

\medskip
\noindent
Failure to select the instructed response would be recorded as a failed
attention check.

\section{Data Processing}

\subsection{Missing, Not-Applicable, and Structurally Unavailable Responses}

\begin{table}[H]
\centering
\caption{Treatment of Missing, Not-Applicable, and Conditional Responses}
\label{tab:missing_responses}

\small

\begin{tabularx}{\textwidth}{
    >{\raggedright\arraybackslash}p{3.2cm}
    >{\raggedright\arraybackslash}p{2.5cm}
    >{\raggedright\arraybackslash}X
    >{\raggedright\arraybackslash}X
}
\toprule
\textbf{Situation} &
\textbf{Classification} &
\textbf{Coding and treatment} &
\textbf{Consequence for analysis} \\
\midrule

No memorable personal experience comes to mind
(\(E_{ig}=0\))
&
Structurally unavailable
&
\(V_{ig}^{\mathrm{E}}\) is missing because the valence question is not displayed.
&
The respondent remains in other analyses. The observation is not used in the H6a specification requiring \(V_{ig}^{\mathrm{E}}\).
\\
\addlinespace

No memorable media information comes to mind
(\(I_{ig}=0\))
&
Structurally unavailable
&
\(V_{ig}^{\mathrm{I}}\) is missing because the valence question is not displayed.
&
The respondent remains in other analyses. The observation is not used in the H6b specification requiring \(V_{ig}^{\mathrm{I}}\).
\\
\addlinespace

Respondent has not worked during the relevant period
&
Not applicable
&
\(C_{ig}^{\mathrm{Work}}\) is coded as missing rather than as no workplace contact. The overall contact measure \(C_{ig}\) is based on \(C_{ig}^{\mathrm{NonWork}}\) only.
&
The respondent remains in the analysis. The unavailable workplace-contact measure does not result in respondent-level exclusion.
\\
\addlinespace

Respondent has worked but reports ``never'' having workplace contact with migrant group \(g\)
&
Valid substantive response
&
\(C_{ig}^{\mathrm{Work}}=1\), corresponding to ``never'' on the contact-frequency scale.
&
The response is included normally in the construction of \(C_{ig}\) and is not treated as missing.
\\
\addlinespace

Respondent selects ``Prefer not to say'' for a displayed question
&
Explicit nonresponse
&
The response is recorded as ``Prefer not to say'' and coded as missing for the corresponding variable in the analysis.
&
The respondent remains in analyses that do not require the respective variable. The observation is unavailable for specifications requiring that variable.
\\
\bottomrule
\end{tabularx}

\vspace{0.5em}

\begin{minipage}{\textwidth}
\footnotesize
\textit{Notes:} All displayed survey questions require a response before 
respondents can proceed. Missing values may arise from explicit nonresponse
or ``Don't know'' responses, questions that are not applicable, or questions
that are not displayed because of conditional survey logic. ``Other''
responses are retained as valid responses and are not automatically treated
as missing.
\end{minipage}

\end{table}

\subsection{Analysis-Sample Availability}

\begin{table}[H]
\centering
\caption{Data Requirements for the Main Analyses}
\label{tab:analysis_sample}

\small

\begin{tabularx}{\textwidth}{
    >{\raggedright\arraybackslash}p{3.5cm}
    >{\raggedright\arraybackslash}X
}
\toprule
\textbf{Analysis} & \textbf{Additional data required} \\
\midrule

H1--H3
&
Valid response for the respective vignette outcome
\\
\addlinespace

H4a--H4b
&
Relevant contact-frequency and experience measures for the migrant group presented in the vignette
\\
\addlinespace

H5
&
Relevant media-exposure and media-tone measures for the migrant group presented in the vignette
\\
\addlinespace

H6a
&
\(E_{ig}=1\) and observed \(V_{ig}^{\mathrm{E}}\)
\\
\addlinespace

H6b
&
\(I_{ig}=1\) and observed \(V_{ig}^{\mathrm{I}}\)
\\
\addlinespace

Adjusted specifications
&
Required respondent-level control variables
\\
\bottomrule
\end{tabularx}

\vspace{0.5em}

\begin{minipage}{\textwidth}
\footnotesize
\textit{Notes:} The number of observations used in each specification may differ depending on the availability of the variables required for the respective analysis. Missing information for a variable not required by a particular specification does not by itself result in respondent-level exclusion.
\end{minipage}

\end{table}